\subsection{User stories}
User stories are a collection of features described from the perspective of the application’s user \cite{cohn2004}. Unlike functional requirements, they provide a better understanding of the needs of non-technical users. This approach ensures that the application will be better aligned with business needs.
\subsubsection{Guest User:}

\begin{enumerate}
	\item As a guest user, I can visit the winery's website to check out its offerings.
	\\Definition of done:
	\begin{itemize}
		\item the winery's “Offer” button is easily accessible on the website,
		\item after clicking the "Offer" button, the winery's offer is displayed.
	\end{itemize}
	\item As a guest user, I can create an account to become a registered user and access personalized features.
	\\Definition of done:
	\begin{itemize}
		\item the "Register" button is clearly visible on the home page and the login form,
		\item after clicking the registration button, the registration form is displayed,
		\item the form requires providing necessary details,
		\item the system validates that the email address is not already in use and that the password meets security requirements,
		\item after submitting the form, a confirmation email is sent to the provided address,
		\item upon clicking the activation link in the email, the account is activated in the database,
		\item the system redirects to a success page and logging in is possible.
	\end{itemize}
	\item As a guest user, I can reset my password to regain access to my account.
	\\Definition of done:
	\begin{itemize}
	\item there is a recovery form with input field for email address,
	\item the system sends an email with a unique reset link,
	\item it is possible to set a new password that meets the security strength requirements after clicking the reset link,
	\item the password hash is updated in the database.
	\end{itemize}
	\item As a guest user, I can log in to my account to access the resources stored there.
	\\Definition of done:
	\begin{itemize}
		\item the login button is easily accessible on the home page,
		\item after clicking the login button, the login form is displayed,
		\item the application requires email address and password to log in,
		\item after entering the correct information, the application redirects to the dashboard,
		\item if incorrect login or password is entered, an error message is displayed prompting to try again.
	\end{itemize}
	\item As a user, I want to be able to open a page with detailed information about the wine batch the bottle comes from by scanning a QR code.
	\\Definition of done:
	\begin{itemize}
		\item the QR code on the bottle is possible to scan using mobile device,
		\item after scanning the code, the relevant page can be opened. It contains detailed information about that particular batch of wine.
	\end{itemize}
\end{enumerate}

\subsubsection{Logged-in user:}

\begin{enumerate}
	\item As a logged-in user I can log out of the system so that my account is secure and no one can use it without my knowledge.
	\\Definition of done:
	\begin{itemize}
		\item after clicking the "log out" button, the application logs the user out of the system, session token is invalidated on the backend,
		\item after logging out, the application redirects to the login page.
	\end{itemize}
	\item As a logged-in user I can change my password so that I can maintain the security of my account.
	\\Definition of done:
	\begin{itemize}
		\item after clicking the "change password" option, the application requires providing current password and the new password. New password (which must meet the security strength criteria) needs to be entered twice for confirmation,
		\item after successfully changing the password, the application displays a confirmation pop-up and asks to log in again using the new password,
		\item the password hash is updated in the database.
	\end{itemize}
	\item As a logged-in user I can change my email address so that I can keep my personal information up to date.
	\\Definition of done:
	\begin{itemize}
		\item after clicking the "change email" option, the application asks to provide new email address and confirm the operation,
		\item after successfully changing the email address, the application displays a confirmation pop-up and asks to log in again using the new email and password,
		\item the email address is updated in the database.
	\end{itemize}
	\item As a logged-in user I can change my phone number so that I can keep my personal information up to date.
	\\Definition of done:
	\begin{itemize}
		\item after clicking the "change phone number" option, the application asks to provide new phone number and confirm the operation,
		\item after successfully changing the phone number, the application displays a confirmation pop-up,
		\item the phone number is updated in the database.
	\end{itemize}
	\item As a logged-in user I can delete my account so that I can remove all my personal data from the system in case I no longer want to use the application.
	\\Definition of done:
	\begin{itemize}
		\item validation prevents deletion if the user is the owner of an organization - in such case, the organization must be deleted first,
		\item after clicking the "delete account" option, the application asks to enter their password and confirm the action,
		\item after successfully deleting the account, the application displays a confirmation pop-up and redirects to the login page,
		\item all personal data is removed from the database, user record is preserved and marked as inactive.
	\end{itemize}
	\item As a logged-in user who is not a part of any organization I can create an organization and become its owner so that I can manage my vineyard and winery operations effectively.
	\\Definition of done:
	\begin{itemize}
		\item validation prevents creation if user already has an organization,
		\item after clicking the "create organization" option, the application asks to provide the organization details and confirm the action,
		\item after successfully creating the organization, the application displays a confirmation pop-up and redirects to the organization management page,
		\item the user is assigned the role of "owner" in the newly created organization,
		\item the organization is added to the database and associated with the user's account.
	\end{itemize}
	\item As a logged-in user I want to join an existing organization so that I can contribute to the management of the vineyard.
	\\Definition of done:
	\begin{itemize}
		\item the user email is visible in user dashboard so that the Vineyard Owner can easily add account to the organization,
		\item after the owner adds the user to the organization, the page needs to be refreshed,
		\item after refreshing the page, the application prompts to confirm joining the organization,
		\item after confirming, the application redirects to the organization dashboard and allows access to the features based on assigned role,
		\item the user is associated with the organization in the database.
	\end{itemize}
\end{enumerate}

\subsubsection{Vineyard Worker:}
\begin{enumerate}
	%===========================================================
	%TASKS
	%===========================================================
	\item As an vineyard worker I would like to be able to view information about completed and planned tasks in order to be able to know what was done and what still needs to be done. \\Definition of done:
	\begin{itemize}
		\item the system displays a list of completed and planned tasks,
		\item the list can be filtered by completion status,
		\item for each task, the list shows its name, due date and completion status,
		\item user can see description of the task when they click on it,
	\end{itemize}
	\item As an vineyard worker I would like to be able to store and edit information about completed and planned tasks in order to be able to view them in the future if needed. \\Definition of done:
	\begin{itemize}
		\item there is button to add new task,
		\item there is a button to edit a task next to each task on the list,
		\item the form for adding/editing a task includes task name,
		\item description, due date, equipment used and field on which task is performed are optional,
	\end{itemize}
	\item As an vineyard worker I would like to be able to delete information about completed and planned tasks in order to be able remove tasks that is not relevant anymore. \\Definition of done:
	\begin{itemize}
		\item there is a button to delete a task,
		\item deleting a task requires confirmation,
	\end{itemize}
	\item As an vineyard worker I would like to be able to mark task as completed or not completed for it to disapear or reapear on list of due tasks. \\Definition of done:
	\begin{itemize}
		\item on task list there is button to mark task as completed or not completed,
	\end{itemize}
	%===========================================================
	%SPRAYING
	%===========================================================
	\item As an vineyard worker I would like to be able to view information about spraying in order to inspect past sprayings. \\Definition of done:
	\begin{itemize}
		\item the system displays a list of all sprayings,
		\item for each spraying, the list shows on which field, what substance was used, in what quantity and when the spraying was performed,
		\item if the spraying has due date, it is highlighted in the list,
	\end{itemize}
	\item As an vineyard worker I would like to be able to store and edit information about spraying in order to be able to view them in the future. \\Definition of done:
	\begin{itemize}
		\item there is button to add new spraying,
		\item there is a button to edit a spraying next to each spraying on the list,
		\item the form for adding/editing a spraying includes substance used, quantity, field on which spraying was performed and selector to choose equipment used for spraying,
		\item spraying date is optional but then worker have to set due date for it,
	\end{itemize}
	\item As an vineyard worker I would like to be able to delete information about spraying in order to remove record added by mistake. \\Definition of done:
	\begin{itemize}
		\item there is a button to delete a spraying,
		\item deleting a spraying requires confirmation,
	\end{itemize}
	%===========================================================
	%PEST AND DISEASE
	%===========================================================
	\item As an vineyard worker I would like to be able to view information about pests and plant diseases on vineyard in order to see past issues. \\Definition of done:
	\begin{itemize}
		\item the system displays a list of all pests and plant diseases records,
		\item for each record, the list shows on which field, what pest or disease was observed and when it was observed,
	\end{itemize}
	\item As an vineyard worker I would like to be able to store and edit information about pests and plant diseases on vineyard in order to be able to view them in the future if needed. \\Definition of done:
	\begin{itemize}
		\item there is button to add new pest or disease record,
		\item there is a button to edit a pest or disease record next to each record on the list,
		\item the form for adding/editing a pest or disease record includes pest or disease name, selector for field on which it was observed and date of observation,
	\end{itemize}
	\item As an vineyard worker I would like to be able to delete information about pests and plant diseases on vineyard in order to remove record added by mistake. \\Definition of done:
	\begin{itemize}
		\item there is a button to delete a pest or disease record,
		\item deleting a pest or disease record requires confirmation,
	\end{itemize}
	%===========================================================
	%HARVEST
	%===========================================================
	\item As an vineyard worker I would like to be able to view information about about amount of grapes harvested on vineyard in order compare it with different harvests. \\Definition of done:
	\begin{itemize}
		\item the system displays a list of all harvest records,
		\item for each record, the list shows on which field, what variety of grapes was harvested, in what quantity and when it was harvested,
	\end{itemize}
	\item As an vineyard worker I would like to be able to store and edit information about about amount of grapes harvested  on vineyard in order to be able to view them in the future if needed. \\Definition of done:
	\begin{itemize}
		\item there is button to add new harvest record,
		\item there is a button to edit a harvest record,
		\item worker needs to enter quantity harvested, field on which harvest was performed and enter date of harvest,
	\end{itemize}
	\item As an vineyard worker I would like to be able to delete information about about amount of grapes harvested on vineyard in order to be remove record added by mistake. \\Definition of done:
	\begin{itemize}
		\item there is a button to delete a harvest record,
		\item deleting a harvest record requires confirmation,
	\end{itemize}
\end{enumerate}
\subsubsection{Oenologist:}

\begin{enumerate}
	%===========================================================
	%MUST
	%===========================================================
	\item As an oenologist, I can view a list of all batches of must, so I can compare specifications with others.
	\\Definition of done:
	\begin{itemize}
		\item the system displays a searchable and filterable list of all must batches,
		\item for each batch, the list shows its unique ID, creation date, location, quantity and batches of grapes used to make it.
	\end{itemize}
	\item As an oenologist, I can create a new must batch, so I can use it later.
	\\Definition of done:
	\begin{itemize}
		\item there is a button allowing the user to start a new must batch,
		\item a form appears in which neccessary details might be selected, such as the piece of equipment where the must batch is to begin, as well as one or more batches of fruit and the quantity from which the batch is to be produced,
		\item there is a button to confirm the creation of a new must batch,
		\item a pop-up appears with confirmation for creating a new must batch,
		\item the system automatically generates a batch number.
	\end{itemize}
	\item As an oenologist, I can update a must batch, so I can correct mistakes I have made during creation.
	\\Definition of done:
	\begin{itemize}
		\item there is a button to edit a must batch,
		\item a form appears where corrects the information can be entered,
		\item a pop-up appears with confirmation for updating a must batch,
		\item the system updates the batch.
	\end{itemize}
	\item As an oenologist, I can delete a must batch, because it doesn't meet my expectations.
	\\Definition of done:
	\begin{itemize}
		\item there is a button to delete a must batch,
		\item a pop-up appears with confirmation for deleting a must batch,
		\item the system correctly deletes the must batch.
	\end{itemize}
	\item As an oenologist, I can scan a QR code on a piece of equipment, so I can quickly view must batch associated with it.
	\\Definition of done:
	\begin{itemize}
		\item QR code located on a vineyard element can be scanned using a mobile device,
		\item the system redirects to the must batch list page,
		\item the system filters the list to show batch related to that specific equipment and highlights the relevant record.
	\end{itemize}
	%===========================================================
	%FERMENTING MUST
	%===========================================================
	\item As an oenologist, I can view a list of all batches of fermenting must, so I can compare specifications with others.
	\\Definition of done:
	\begin{itemize}
		\item the system displays a searchable and filterable list of all fermenting must batches,
		\item for each batch, the list shows its unique ID, creation date, location, quantity and batches of must used to make it.
	\end{itemize}
	\item As an oenologist, I can create a new fermenting must batch, so I can use it later.
	\\Definition of done:
	\begin{itemize}
		\item there is a button to start a new fermenting must batch,
		\item a form appears where details can be entered, such as the piece of equipment where the fermenting must batch is to begin, as well as one or more batches of must and the quantity from which the batch is to be produced,
		\item there is a button to confirm the creation of a new fermenting must batch,
		\item a pop-up appears with confirmation for creating a new fermenting must batch,
		\item the system automatically generates a batch number.
	\end{itemize}
	\item As an oenologist, I can update a fermenting must batch, so I can correct mistakes I have made during creation.
	\\Definition of done:
	\begin{itemize}
		\item there is a button to edit a fermenting must batch,
		\item a form appears where the information can be corrected,
		\item a pop-up appears with confirmation for updating a fermenting must batch,
		\item the system updates the batch.
	\end{itemize}
	\item As an oenologist, I can delete a fermenting must batch, because it doesn't meet my expectations.
	\\Definition of done:
	\begin{itemize}
		\item there is a button to delete a fermenting must batch,
		\item a pop-up appears with confirmation for deleting a fermenting must batch,
		\item the system correctly deletes the fermenting must batch.
	\end{itemize}
	%===========================================================
	%INGREDIENTS
	%===========================================================
	\item As an oenologist, I can view a list of all ingredients used during production of a specific batch, so I can use this information in the future.
	\\Definition of done:
	\begin{itemize}
		\item the system displays a list of ingredients used during production of a specific batch.
	\end{itemize}
	\item As an oenologist, I can add ingredients and their quantities to log this data for future use.
	\\Definition of done:
	\begin{itemize}
		\item there is a button to add ingredients to a specific production batch of must,
		\item a form appears where the ingredients and their quantity can be entered,
		\item there is a button to confirm the addition of a new ingredient to a batch,
		\item a pop-up appears with confirmation for adding a new ingredient,
		\item the system automatically generates an ingredient number.
	\end{itemize}
	\item As an oenologist, I can update an ingredient record, so I can correct mistakes I have made during creation.
	\\Definition of done:
	\begin{itemize}
		\item there is a button to edit an ingredient record,
		\item a form appears where the information can be corrected,
		\item a pop-up appears with confirmation for updating an ingredient record,
		\item the system updates the ingredient record.
	\end{itemize}
	\item As an oenologist, I can delete an ingredient, because I added it to the wrong production batch of must.
	\\Definition of done:
	\begin{itemize}
		\item there is a button to delete an ingredient,
		\item a pop-up appears with confirmation for deleting an ingredient,
		\item the system correctly deletes the ingredient.
	\end{itemize}
	%===========================================================
	%AGED WINE
	%===========================================================
	\item As an oenologist, I can view a list of all batches of aged wine, so I can compare specifications with others.
	\\Definition of done:
	\begin{itemize}
		\item the system displays a searchable and filterable list of all aged wine batches,
		\item for each batch, the list shows its unique ID, creation date, location, quantity and batches of fermented must used to make it.
	\end{itemize}
	\item As an oenologist, I can create a new aged wine batch, so I can use it later.
	\\Definition of done:
	\begin{itemize}
		\item there is a button to start a new aged wine batch,
		\item a form appears where details can be selected, such as the piece of equipment where the aged wine batch is to begin, as well as one or more batches of fermented must and the quantity from which the batch is to be produced,
		\item there is a button to confirm the creation of a new aged wine batch,
		\item a pop-up appears with confirmation for creating a new aged wine batch,
		\item the system automatically generates a batch number.
	\end{itemize}
	\item As an oenologist, I can update an aged wine batch, so I can correct mistakes I have made during creation.
	\\Definition of done:
	\begin{itemize}
		\item there is a button to edit an aged wine batch,
		\item a form appears where the information can be corrected,
		\item a pop-up appears with confirmation for updating an aged wine batch,
		\item the system updates the batch.
	\end{itemize}
	\item As an oenologist, I can delete an aged wine batch, because it doesn't meet my expectations.
	\\Definition of done:
	\begin{itemize}
		\item there is a button to delete an aged wine batch,
		\item a pop-up appears with confirmation for deleting an aged wine batch,
		\item the system correctly deletes the aged wine batch.
	\end{itemize}
	%===========================================================
	%BLENDED WINE
	%===========================================================
	\item As an oenologist, I can view a list of all batches of blended wine, so I can compare specifications with others.
	\\Definition of done:
	\begin{itemize}
		\item the system displays a searchable and filterable list of all blended wine batches,
		\item for each batch, the list shows its unique ID, creation date, location, quantity and batches of wine (fermented must) and/or aged wine used to make it.
	\end{itemize}
	\item As an oenologist, I can create a new blended wine batch, so I can use it later.
	\\Definition of done:
	\begin{itemize}
		\item there is a button to start a new blended wine batch,
		\item a form appears where details can be selected, such as the piece of equipment where the blended wine batch is to begin, as well as one or more batches of wine (fermented must) and/or aged wine and the quantity from which the batch is to be produced,
		\item there is a button to confirm the creation of a new blended wine batch,
		\item a pop-up appears with confirmation for creating a new blended wine batch,
		\item the system automatically generates a batch number.
	\end{itemize}
	\item As an oenologist, I can update a blended wine batch, so I can correct mistakes I have made during creation.
	\\Definition of done:
	\begin{itemize}
		\item there is a button to edit a blended wine batch,
		\item a form appears where the information can be corrected,
		\item a pop-up appears with confirmation for updating a blended wine batch,
		\item the system updates the batch.
	\end{itemize}
	\item As an oenologist, I can delete a blended wine batch, because it doesn't meet my expectations.
	\\Definition of done:
	\begin{itemize}
		\item there is a button to delete a blended wine batch,
		\item a pop-up appears with confirmation for deleting a blended wine batch,
		\item the system correctly deletes the blended wine batch.
	\end{itemize}
	%===========================================================
	%BOTTLED WINE
	%===========================================================
	\item As an oenologist, I can view a list of all batches of bottled wine, so I can compare specifications with others.
	\\Definition of done:
	\begin{itemize}
		\item the system displays a searchable and filterable list of all bottled wine batches,
		\item for each batch, the list shows its unique ID, creation date, location, quantity and batches of wine (fermented must) and/or aged wine and/or blended wine used to make it.
	\end{itemize}
	\item As an oenologist, I can create a new bottled wine batch, so I can end the production.
	\\Definition of done:
	\begin{itemize}
		\item there is a button to start a new bottled wine batch,
		\item a form appears where details can be selected, such as the piece of equipment where the bottled wine batch is to begin, as well as one or more batches of wine (fermented must) and/or aged wine and/or blended wine and the quantity from which the batch is to be produced,
		\item there is a button to confirm the creation of a new bottled wine batch,
		\item a pop-up appears with confirmation for creating a new bottled wine batch,
		\item the system automatically generates a batch number.
	\end{itemize}
	\item As an oenologist, I can update a bottled wine batch, so I can correct mistakes I have made during creation.
	\\Definition of done:
	\begin{itemize}
		\item there is a button to edit a bottled wine batch,
		\item a form appears where the information can be corrected,
		\item a pop-up appears with confirmation for updating a bottled wine batch,
		\item the system updates the batch.
	\end{itemize}
	\item As an oenologist, I can delete a bottled wine batch, because it doesn't meet my expectations.
	\\Definition of done:
	\begin{itemize}
		\item there is a button to delete a bottled wine batch,
		\item a pop-up appears with confirmation for deleting a bottled wine batch,
		\item the system correctly deletes the bottled wine batch.
	\end{itemize}
	%===========================================================
	%Ingredient
	%===========================================================
	\item As an oenologist, I can view a list of all ingredients, so I can see what I have already added.
	\\Definition of done:
	\begin{itemize}
		\item the system displays a searchable and filterable list of all ingredients.
	\end{itemize}
	\item As an oenologist, I can add a new ingredient, so I can use it later.
	\\Definition of done:
	\begin{itemize}
		\item there is a button to add a new ingredient,
		\item a form appears in which a new ingredient can be created,
		\item there is a button to confirm the addition of a new ingredient,
		\item a pop-up appears with confirmation for adding a new ingredient,
		\item the system automatically generates an ingredient number.
	\end{itemize}
	\item As an oenologist, I can update a ingredient, so I can correct mistakes I have made during creation.
	\\Definition of done:
	\begin{itemize}
		\item there is a button to edit a ingredient,
		\item a form appears where the information can be corrected,
		\item a pop-up appears with confirmation for updating a ingredient,
		\item the system updates the ingredient.
	\end{itemize}
	\item As an oenologist, I can delete a ingredient, because it doesn't meet my expectations.
	\\Definition of done:
	\begin{itemize}
		\item there is a button to delete a must batch,
		\item a pop-up appears with confirmation for deleting a ingredient,
		\item the system correctly deletes the ingredient.
	\end{itemize}
	%===========================================================
	% CHEMICAL PARAMETERS
	%===========================================================
	\item As an oenologist, I can record and view chemical parameters for any production batch, to monitor quality.
	\\Definition of done:
	\begin{itemize}
		\item the batch can be selected (manually or via QR code) and parameters can be entered into a form,
		\item the form includes sugar level, acidity, pH, temperature, and alcohol content,
		\item the system stores these values with a timestamp and links them to the specific batch.
	\end{itemize}
	%===========================================================
	% SPOILED STATUS MANAGEMENT 
	%===========================================================
	\item As an oenologist, I can mark or unmark a batch as spoiled, to prevent it from entering the next production stage.
	\\Definition of done:
	\begin{itemize}
		\item the batch can be selected (manually or via QR code),
		\item there is a toggle for the "Spoiled" status,
		\item a confirmation pop-up appears,
		\item the system visually highlights spoiled batches in all lists and prevents further processing unless unmarked.
	\end{itemize}
	%===========================================================
	% QR SCANNING GENERAL
	%===========================================================
	\item As an oenologist, I can use QR codes to quickly access and interact with any production stage.
	\\Definition of done:
	\begin{itemize}
		\item scanning a QR code on winery equipment redirects the user to the correct category (must, fermenting, etc.),
		\item the system automatically filters for the specific batch in that location and highlights it,
		\item if the equipment is empty, the system suggests creating a new batch for that location.
	\end{itemize}
	%===========================================================
	% potassium metabisulfite
	%===========================================================
	\item As an oenologist, I can see how much, potassium metabisulfite, I should add in order to stop the fermentation process, so I do not need to calculate it by hand.
	\\Definition of done:
	\begin{itemize}
		\item the batch can be selected (manually or via QR code),
		\item the calculate potassium metabisulfite button is highlighted and easily accessible,
		\item the system automatically calculates the required amount of potassium metabisulfite based on the current quantity and specific parameters,
		\item the system displays the calculated dosage,
		\item the system displays details, such as the formula and change the amount of potassium metabisulfite if needed.
	\end{itemize}
	%===========================================================
	% predicted alcohol level
	%===========================================================
	\item As an oenologist, I can see the predicted alcohol level in wine after fermentation, so that I will know how much sugar do I need to add to reach the desired alcohol content.
	\\Definition of done:
	\begin{itemize}
		\item the batch can be selected (manually or via QR code),
		\item the system displays the predicted alcohol content for a batch of juice or a fermented batch, taking into account the sugar content from the fruit and any added sugar, as well as the type of yeast used.
	\end{itemize}
	%===========================================================
	% nutrients
	%===========================================================
	\item As an oenologist, I can see the calculated nutrients, so I will know, how much of every nutrient is in the final product.
	\\Definition of done:
	\begin{itemize}
		\item after finalizing production, the show nutrients button is displayed for every finalized batch,
		\item after clicking on the show nutrients button, the table of nutrients is displayed,
		\item the nutrients parameters can be changed if needed.
	\end{itemize}
\end{enumerate}

\subsubsection{Vineyard Owner:}

\begin{enumerate}
	%===========================================================
	% REPORTS
	%===========================================================
	\item As a vineyard owner, I can generate production efficiency report in order to review it and analyze the efficiency of the production process.
	\\Definition of done:
	\begin{itemize}
		\item after selecting the option to generate production efficiency report, details can be selected, such as the time period and/or specific production batches for which the report is to be generated,
		\item after confirming the selected parameters, a report is generated,
		\item the report includes key performance indicators such as yield per hectare, average time from harvest to bottling, and percentage of batches marked as spoiled,
		\item the report is saved in the system and can be accessed later from the reports section.
	\end{itemize}
	
	\item As a vineyard owner, I can generate diversity of grapevine varieties report in order to review it and analyze the variety distribution in the vineyard.
	\\Definition of done:
	\begin{itemize}
		\item after selecting the option to generate diversity of grapevine varieties report, details can be selected, such as the time period and/or specific fields for which the report is to be generated,
		\item after confirming the selected parameters, a report is generated,
		\item the report includes information about the different grapevine varieties planted in the vineyard, their respective areas, and their contribution to the overall production,
		\item the report is saved in the system and can be accessed later from the reports section.
	\end{itemize}
	
	\item As a vineyard owner, I can generate vineyard operations report in order to review it and analyze the operational efficiency of the vineyard.
	\\Definition of done:
	\begin{itemize}
		\item after selecting the option to generate vineyard operations report, details can be selected, such as the time period and/or specific operations for which the report is to be generated,
		\item after confirming the selected parameters, a report is generated,
		\item the report includes information about the various operations performed in the vineyard (e.g., pruning, irrigation, pest control), their frequency, and their impact on production efficiency,
		\item the report is saved in the system and can be accessed later from the reports section.
	\end{itemize}
	
	\item As a vineyard owner, I can generate historical trend report in order to review it and analyze the production trends over time.
	\\Definition of done:
	\begin{itemize}
		\item after selecting the option to generate historical trend report, details can be selected, such as the time period, and aggregation level (e.g., monthly, quarterly, yearly),
		\item after confirming the selected parameters, a report is generated,
		\item the report includes the same information as other reports, but aggregated over time to show trends and patterns,
		\item the report is saved in the system and can be accessed later from the reports section.
	\end{itemize}
	
	%===========================================================
	% ORGANIZATION
	%===========================================================
	\item As a vineyard owner, I can view information about the organization.
	\\Definition of done:
	\begin{itemize}
		\item the information includes the organization name, address, contact information, and a list of associated vineyards and winery equipment.
	\end{itemize}
	\item As a vineyard owner, I can view information about the organization.
	\\Definition of done:
	\begin{itemize}
		\item the information includes the organization name, address, contact information, and a list of associated vineyards and winery equipment.
	\end{itemize}
	
	\item As a vineyard owner, I can edit information about the organization, in order to keep it up-to-date.
	\\Definition of done:
	\begin{itemize}
		\item every field of the organization information can be modified, including the name, address and company logo,
	\end{itemize}
	
	\item As a vineyard owner, I can remove the organization from the system, which will also remove all associated data, in order to ensure that organization data is not stored in the system if I no longer want to use the application.
	\\Definition of done:
	\begin{itemize}
		\item the removal action must be confirmed by entering the password,
		\item upon confirmation, the organization and all associated data apart from wine bottled in the past is deleted from the system,
		\item after successful deletion, the vineyard owner role is revoked, and the user is redirected to the home page.
	\end{itemize}
	
	%===========================================================
	% FIELDS
	%===========================================================
	\item As a vineyard owner, I can view information about fields from the organization along with the grape variety planted there, in order to see what type of grapes I grow.
	\\Definition of done:
	\begin{itemize}
		\item there is a list of all fields from the organization,
		\item for each field, the list shows its area and what grape variety is planted there,
	\end{itemize}
	\item As a vineyard owner, I can add and update information about vineyard fields from the organization along with the grape variety planted there, in order to keep the information about the vineyard up-to-date.
	\\Definition of done:
	\begin{itemize}
		\item there is a button to add a new field,
		\item there is a button to edit an existing field next to each field on the list,
		\item the form for adding/editing a field includes fields for the area and selector for grape variety,
	\end{itemize}
	\item As a vineyard owner, I can delete information about vineyard fields from the organization along with the grape variety planted there, in order to be able to remove fields that do not belong to me anymore.
	\\Definition of done:
	\begin{itemize}
		\item there is a button to delete a field,
		\item deleting a field requires confirmation,
		\item the field should still be linked to the historical data, but it should be marked as "unactive" and not appear in the list of available fields.
	\end{itemize}
	
	%===========================================================
	% SUBSTANCES
	%===========================================================
	\item As a vineyard owner, I can view information about substances used in the vineyard, in order to be hand it in in case of enviornment control.
	\\Definition of done:
		\begin{itemize}
		\item there is a list of all substances,
		\item for each substance, the list shows its name,
	\end{itemize}
	\item As a vineyard owner, I can store and edit information about substances used in the vineyard, in order to keep the information about substances up-to-date.
	\\Definition of done:
	\begin{itemize}
		\item there is a button to add a new substance,
		\item there is a button to edit an existing substance next to each substance on the list,
		\item the form for adding/editing a substance includes field for the name,
	\end{itemize}
	\item As a vineyard owner, I can delete information about substances used in the vineyard, in order to be able to remove substances that I do not use anymore.
	\\Definition of done:
	\begin{itemize}
		\item there is a button to delete a substance,
		\item deleting a substance requires confirmation,
		\item the substance should still be linked to the historical data, but it should be marked as "unactive" and not appear in the list of available substances.
	\end{itemize}
	%===========================================================
	% EQUIPMENT
	%===========================================================
	\item As a vineyard owner, I can store and edit information about equipment, in order to link it to batches, tasks etc.
	\\Definition of done:
	\begin{itemize}
		\item there is a button to add a new piece of equipment,
		\item there is a button to edit an existing piece of equipment next to each piece of equipment on the list,
		\item the form for adding/editing a piece of equipment includes fields for the name and type selector,
		\item there is option to note date of end of technical inspection,
	\end{itemize}
	\item As a vineyard owner, I can view information about equipment, in order to see what equipment I own.
	\\Definition of done:
	\begin{itemize}
		\item there is a list of all equipment,
		\item for each piece of equipment, the list shows its name, type and date of end of technical inspection (if noted),
	\end{itemize}
	\item As a vineyard owner, I can delete information about equipment, in order to remove data about equipment that is out of use.
	\\Definition of done:
	\begin{itemize}
		\item there is a button to delete a piece of equipment,
		\item deleting a piece of equipment requires confirmation,
		\item the equipment should still be linked to the historical data, but it should be marked as "unactive" and not appear in the list of available equipment.
	\end{itemize}

	%===========================================================
	% OFFERINGS
	%===========================================================
	\item As a vineyard owner, I can add and edit offerings from front page, in order for customers to be able to read information about wine.
	\\Definition of done:
	\begin{itemize}
		\item there is a button to add a new offering,
		\item there is a button to edit an existing offering next to each offering on the list,
		\item the form for adding/editing an offering includes selectors for the batch of wine and field for price,
	\end{itemize}
	\item As a vineyard owner, I can view offerings from front page, in order to see what wines I have on offer.
	\\Definition of done:
	\begin{itemize}
		\item there is a list of all offerings on the front page,
		\item for each offering, the list shows the name of the wine and its price,
	\end{itemize}
	\item As a vineyard owner, I can delete offerings from front page, in order to remove wine from offer.
	\\Definition of done:
	\begin{itemize}
		\item there is a button to delete an offering,
		\item deleting an offering requires confirmation,
	\end{itemize}

	%===========================================================
	% USERS
	%===========================================================
	\item As a vineyard owner, I can add users to my organization, in order to allow them to access the system and contribute to the management of the vineyard.
	\\Definition of done:
	\begin{itemize}
		\item the vineyard owner is able to enter the email address provided by the employee they want to add to the organization, 
		\item the system checks if there is an existing account associated with that email address,
		\item if there is an existing account, the user is added to the organization,
		\item after employee confirms joining the organization, new user appears in the user list,
		\item if there is no existing account, error message is displayed, prompting the vineyard owner to ask the employee to register first.
	\end{itemize}
	
	\item As a vineyard owner, I can view users in my organization, in order to manage the access to the system and ensure that only authorized personnel can access the vineyard data.
	\\Definition of done:
	\begin{itemize}
		\item there is a list of all users in the organization,
		\item users' information can be edited,
		\item there is a button for removing users from the organization,
	\end{itemize}
	\item As a vineyard owner, I can delete users from organization, in order to remove account of people who don't work for me anymore.
	\\Definition of done:
	\begin{itemize}
		\item there is a button for removing users from the organization,
		\item deleting a user requires confirmation,
	\end{itemize}
	
	\item As a vineyard owner, I can assign and revoke roles to other users in the organization, in order to manage their permissions and access to different functionalities of the system.
	\\Definition of done:
	\begin{itemize}
		\item the roles can be assigned to users, such as vineyard worker or oenologist,
		\item multiple roles can be assigned to a single user,
		\item the roles can be revoked from users,
	\end{itemize}
\end{enumerate}
